\documentclass[11pt,letterpaper]{article}
\usepackage{iftex}
\ifPDFTeX\usepackage[utf8]{inputenc}\fi
\usepackage[T1]{fontenc}
\usepackage[spanish,es-nodecimaldot]{babel}
\usepackage[margin=2.2cm,headheight=15pt]{geometry}
\usepackage{lmodern,microtype,amsmath}
\usepackage[table]{xcolor}
\usepackage{array,booktabs,tabularx,enumitem,fancyhdr,hyperref}
\definecolor{marino}{HTML}{002D62}
\definecolor{dorado}{HTML}{F2A900}
\hypersetup{colorlinks=true,linkcolor=marino,urlcolor=marino}
\pagestyle{fancy}
\fancyhf{}
\fancyhead[L]{\small EmprendeBúho}
\fancyhead[R]{\small Planeación, seguimiento y control}
\fancyfoot[C]{\thepage}
\setlength{\parindent}{0pt}
\setlength{\parskip}{6pt}
\setlength{\emergencystretch}{2em}
\renewcommand{\arraystretch}{1.16}
\setlist{nosep,leftmargin=*}
\newcolumntype{Y}{>{\raggedright\arraybackslash}X}
\newcommand{\hu}[1]{\mbox{HU-#1}}

\begin{document}
\begin{center}
{\Large\bfseries Universidad de Sonora}\par
{\large Ingeniería de Software II}\par
\vspace{0.3cm}
{\LARGE\bfseries\color{marino} Evaluación de planeación,\\
seguimiento y control de proyectos}\par
\vspace{0.25cm}
{\large\bfseries Proyecto EmprendeBúho}\par
Equipo \guillemotleft Descompuestos por Cholesky\guillemotright\par
Profesor: Adrian Vazquez Osorio\par
Hermosillo, Sonora, 10 de noviembre de 2026
\end{center}
\hrule

\textbf{Base de las respuestas.} Se utiliza el \textit{Avance de Proyecto
(Parte 01)} para historias, prioridades y red de tareas; las bitácoras
originales para hechos documentados; y el último estado comunicado por el
equipo para el tablero: 36 historias terminadas, una en proceso y tres
pendientes. Las estimaciones nuevas y los ejemplos hipotéticos se identifican
expresamente. El equipo confirmó 27 revisiones informales al 10 de noviembre
y el cierre real de T1 el 19 de octubre. Estas confirmaciones actualizan los
datos que antes estaban pendientes; las demás entradas reconstruidas siguen
sujetas a validación.

\section*{1. Historia de usuario prioritaria (2 puntos)}
\textbf{Seleccionamos la \hu{01}: solicitud de cuenta del emprendedor.}

\begin{quote}
Como emprendedor de la Universidad de Sonora, quiero solicitar una cuenta
para registrar mi negocio en la plataforma.
\end{quote}

En el primer avance tiene \textbf{prioridad alta y complejidad alta}. La
consideramos prioritaria porque habilita la incorporación del negocio y su
posterior revisión; sin cuentas de emprendedores no puede alimentarse el
directorio. El documento contiene otras historias de prioridad alta: esta
selección se justifica por su papel habilitador, sin suponer que existe una
clasificación original que la coloque por encima de todas las demás.

\section*{2. Historia más corta y duración estimada (2 puntos)}
\textbf{Elegimos la \hu{05}: clasificar el negocio por categoría.}

\begin{quote}
Como emprendedor, quiero asignar una categoría a mi negocio para facilitar
que los clientes lo encuentren.
\end{quote}

\textbf{Duración estimada para esta evaluación: 1 día hábil.} Su complejidad
original es baja. El alcance comprende seleccionar una categoría existente,
guardar la asociación y revisar que se muestre correctamente. Se supone que
ya existen la base de datos, las categorías y el formulario del perfil; no
incluye desarrollar el módulo administrativo de categorías.

La Parte 01 asigna días a tareas agrupadas, no a historias individuales. Por
eso, un día es una \textbf{estimación propuesta}, no una duración registrada
allí. Otras historias de complejidad baja pueden tener una duración similar;
se elige HU-05 como una de las más cortas bajo este alcance.

\section*{3. Historias en la columna Por hacer (2 puntos)}
\textbf{Quedan 3 historias en Por hacer:}
\begin{itemize}
\item \hu{03}: acceso de clientes al sitio web sin registro, con el alcance
corregido que sustituye la antigua foto de menú.
\item \hu{10}: publicar la ubicación aproximada del negocio.
\item \hu{23}: buscar emprendimientos por zona del campus.
\end{itemize}

La \hu{30}, ordenamiento de resultados, está \textbf{en proceso}, por lo que
no se cuenta en Por hacer. La comprobación del total es:
\[
36\text{ terminadas}+1\text{ en proceso}+3\text{ por hacer}
=\boxed{40\text{ historias}}.
\]
El conteo procede del estado comunicado por el equipo, no de una consulta en
vivo a Trello. La modificación de HU-03 no aumenta el total de historias.

\section*{4. Dos rutas dependientes y sus duraciones (2 puntos)}
Para este cálculo utilizamos \textbf{las duraciones estimadas de la tabla de
dependencias de la Parte 01}, en días hábiles de trabajo:

\begin{center}\small
\begin{tabularx}{\linewidth}{@{}lYrY@{}}
\toprule
Tarea & Actividad & Días & Predecesoras \\\midrule
T1 & Arquitectura y base de datos & 5 & Ninguna \\
T2 & Gestión de cuentas y autenticación & 10 & T1 \\
T3 & Perfil del emprendimiento & 10 & T2 \\
T4 & Catálogo y productos & 10 & T3 \\
T5 & Visualización, búsqueda y filtros & 15 & T4 y T8 \\
T6 & Interacción del consumidor & 5 & T5 \\
T7 & Administración & 10 & T2 \\
T8 & Ubicación del emprendimiento & 8 & T3 \\
T9 & Pruebas y entrega & 7 & T6 y T7 \\\bottomrule
\end{tabularx}
\end{center}

\textbf{Ruta A:} T1 $\rightarrow$ T2 $\rightarrow$ T3 $\rightarrow$ T4
$\rightarrow$ T5 $\rightarrow$ T6 $\rightarrow$ T9.
\[
D_A=5+10+10+10+15+5+7=\boxed{62\text{ días}}.
\]

\textbf{Ruta B:} T1 $\rightarrow$ T2 $\rightarrow$ T3 $\rightarrow$ T8
$\rightarrow$ T5 $\rightarrow$ T6 $\rightarrow$ T9.
\[
D_B=5+10+10+8+15+5+7=\boxed{60\text{ días}}.
\]

T5 requiere tanto T4 como T8. La suma de cada ruta describe una cadena de
dependencias; T4 y T8 pueden realizarse en paralelo y no se suman una detrás
de la otra.

\section*{5. Ruta crítica y holgura (2 puntos)}
\textbf{La ruta crítica es la ruta A, con 62 días.} La otra ruta completa
de la red, T1 $\rightarrow$ T2 $\rightarrow$ T7 $\rightarrow$ T9, suma
$5+10+10+7=32$ días, por lo que tampoco supera a A. Se supone la red de
precedencias original, sin restricciones adicionales de recursos.

La diferencia entre las dos rutas elegidas es:
\[
H=D_A-D_B=62-60=\boxed{2\text{ días hábiles}}.
\]

T4 y T8 comienzan después de T3: T4 dura 10 días y T8 dura 8. Por ello, T8
puede demorarse hasta 2 días sin retrasar el inicio de T5, suponiendo que las
demás duraciones se mantienen. Esta es la holgura de la rama alternativa;
\textbf{las actividades compartidas de la ruta crítica no adquieren esa
holgura}.

\textbf{Consistencia del cálculo.} Las fechas del Gantt abarcan 45 días
hábiles, pero no coinciden con todas las duraciones de su tabla de
estimaciones. Los resultados 62, 60 y 2 pertenecen a esa tabla original;
no describen las ventanas superpuestas del cronograma de recuperación de la
Parte 02. No deben combinarse ambos modelos como si fueran uno solo.

\section*{6. Probabilidad, impacto y puntuación de un riesgo (2 puntos)}
\textbf{Riesgo:} retraso en la publicación de ubicación (HU-10), que impida
cerrar la búsqueda por zona (HU-23) y la cercanía del ordenamiento (HU-30).

Para esta evaluación asignamos los siguientes valores de juicio del equipo;
no son frecuencias calculadas a partir de una serie histórica:
\begin{itemize}
\item Probabilidad: \textbf{75\,\%}, equivalente a $0.75$.
\item Impacto: \textbf{3 en una escala de 1 a 4}, por afectar varias funciones y su integración.
\end{itemize}
\[
\text{Puntuación}=P\times I=0.75\times3=\boxed{2.25}.
\]
El resultado es una puntuación de exposición al riesgo, no un porcentaje ni
un número de días de retraso. Como respuesta, Orlando, Ángel y Leinad deben
acordar el contrato de zonas, revisar los bloqueos diariamente e integrar
primero HU-10, después HU-23 y finalmente la aceptación de cercanía de HU-30.

\section*{7. Revisiones informales realizadas (2 puntos)}
\textbf{Hemos realizado 27 revisiones informales de avance}, según el conteo
confirmado por el equipo al 10 de noviembre de 2026.

Para este seguimiento, cada registro diario de la bitácora del equipo cuenta
como una revisión informal: se revisan las actividades anteriores, el trabajo
del día y las necesidades o bloqueos de los integrantes. El cálculo es:
\[
5\text{ semanas}\times5\text{ revisiones por semana}
+2\text{ días de la semana 6}=\boxed{27\text{ revisiones}}.
\]
Las dos revisiones de la semana 6 corresponden al 9 y 10 de noviembre. Se
cuenta una revisión por jornada del equipo, no una por integrante; por ello
no se multiplica por los ocho participantes. Tampoco se cuentan los días
11--13 de noviembre, que todavía están previstos en el borrador semanal.

\section*{8. Diseño hipotético de una revisión formal (4 puntos)}
\textbf{Entregable elegido: módulo de gestión de cuentas y autenticación
(T2)}, correspondiente a HU-01, HU-27, HU-28, HU-32 y HU-40 de la Parte 01.
La siguiente es una propuesta de revisión, no una reunión ya realizada.

\textbf{Objetivo:} determinar si la versión entregada cumple los criterios
acordados de solicitud de cuenta, acceso por rol, notificaciones, recuperación
de contraseña y baja de cuenta, y si puede aceptarse para la entrega final.

\textbf{Participantes propuestos:} Denzel como moderador; Fernanda y Mario
como representantes de los autores del módulo; Ángel y Orlando como
revisores; y Leinad como secretario. Los autores de una parte específica
explican su trabajo, pero no actúan como únicos aprobadores de esa parte.

\textbf{Proceso de revisión formal:}
\begin{enumerate}
\item \textbf{Planear y convocar.} Definir fecha, duración, alcance,
participantes y criterios de entrada. Identificar una versión fija del
entregable y distribuirla junto con los requisitos y la lista de verificación.
Los participantes confirman que el material está completo antes de reunirse.
\item \textbf{Preparación individual.} Cada revisor inspecciona los flujos,
formularios, validaciones y evidencia de pruebas. Registra por separado
defectos, dudas y el requisito afectado, para no depender de una demostración
improvisada durante la sesión.
\item \textbf{Inspección moderada.} Recorrer escenarios preparados: solicitud
válida e inválida; cuenta pendiente, aprobada, rechazada y suspendida;
credenciales incorrectas; código de recuperación vencido; eliminación de
cuenta; y restricción de rutas por rol. Verificar que las credenciales
estudiantiles no se expongan públicamente.
\item \textbf{Registrar hallazgos.} El secretario elabora el acta y la lista
de defectos con identificador, descripción, severidad, evidencia,
responsable y fecha de corrección. La sesión se centra en identificar
problemas; las correcciones se realizan después.
\item \textbf{Emitir un dictamen.} Aceptar, aceptar con correcciones menores
o rechazar para retrabajo, según criterios definidos previamente. No se
acepta el módulo si hay fallos críticos de acceso, privacidad o pérdida de
datos, o si faltan criterios obligatorios.
\item \textbf{Corregir y verificar.} Los autores atienden los hallazgos.
Los revisores comprueban las correcciones y revisan que no introduzcan
regresiones. Si el cambio es amplio, se programa otra sesión formal.
\item \textbf{Cerrar y conservar evidencia.} El moderador registra el
dictamen final, la versión aceptada y la aprobación de los participantes.
Se archivan el acta, la lista de verificación, los resultados de pruebas
y el seguimiento de defectos; después se actualiza el tablero.
\end{enumerate}

\textbf{Productos de la revisión:} acta formal, lista de defectos, acuerdos
con responsables y fechas, evidencia de verificación y dictamen de aceptación.
Estos registros permiten distinguirla de una conversación informal de avance.

\section*{9. Desviación entre días planificados y reales (2 puntos)}
\textbf{La desviación de plazo de T1 es de +11 días hábiles de retraso.}
El equipo confirmó que arquitectura y base de datos terminaron el
\textbf{19 de octubre de 2026}; el cierre planificado era el
\textbf{2 de octubre de 2026}.

El Gantt original fijó T1 del 28 de septiembre al 2 de octubre:
\textbf{5 días hábiles}. Desde ese mismo inicio planificado hasta el cierre
real del 19 de octubre transcurrieron \textbf{16 días hábiles}, contando
ambos extremos, lunes a viernes y sin descontar feriados.
\[
\Delta_{T1}=16-5=\boxed{+11\text{ días hábiles}}.
\]
La misma diferencia se obtiene contando las once jornadas hábiles posteriores
al 2 de octubre y hasta el 19 de octubre, inclusive. El signo positivo
indica retraso respecto al plazo del hito. Se cuenta desde el
\textbf{inicio planificado} para incluir el retraso de arranque; dieciséis
días transcurridos no equivalen a dieciséis jornadas efectivas trabajadas.
Esta es la desviación de T1, no una demora final demostrada de todo el proyecto.

Para el proyecto completo, la línea base comprende 45 días hábiles del
28 de septiembre al 27 de noviembre. Al 10 de noviembre han transcurrido
32 jornadas de ese calendario. Restar $32-45=-13$ \textbf{no demostraría un
adelanto}: el proyecto sigue abierto y quedan días de su ventana prevista.
La desviación final se calculará con la fecha real de entrega confirmada.
La fecha del 19 de octubre sustituye al 20 de octubre que figuraba como
propuesta de cierre de T1 en el borrador de la Parte 02.

\section*{Referencias y confirmaciones del equipo}
\begin{itemize}
\item \textit{Copia de AvanceDeProyecto (Parte 01)}, 25 de septiembre de
2026: matriz de historias, prioridad, complejidad, tareas y dependencias.
\item \textit{Bitácoras originales de las semanas 1 y 2}: investigación,
seguimiento, reunión del 13 de octubre y retraso de arquitectura.
\item Descripción del Gantt original: fechas del 28 de septiembre al
27 de noviembre y diferencias respecto a las duraciones de la tabla.
\item Estado comunicado por el equipo: 36 terminadas, HU-30 en proceso y
HU-03, HU-10 y HU-23 pendientes; HU-03 corregida a acceso sin registro.
\end{itemize}
\textbf{Confirmaciones incorporadas:} 27 revisiones informales, calculadas
como cinco semanas de cinco registros diarios más dos días de la semana 6;
y cierre real de T1 el 19 de octubre de 2026. El conteo de revisiones se
atribuye al equipo, sin suponer que todos sus registros estén transcritos
en los PDF originales aportados.

\end{document}
